% ==============================================================================
%  Flow-free topological charge density — design document
%
%  Companion to reports/glueball/glueball_spectroscopy.tex and
%  reports/attention/attention_as_operator.tex; self-contained.
%  Build: pdflatex (run twice for refs).
% ==============================================================================
\documentclass[11pt,a4paper]{article}

\usepackage[margin=2.6cm]{geometry}
\usepackage[T1]{fontenc}
\usepackage[utf8]{inputenc}
\usepackage{lmodern}
\usepackage{microtype}
\usepackage{amsmath,amssymb,mathtools}
\usepackage{graphicx}
\usepackage{booktabs}
\usepackage{caption}
\usepackage{xcolor}
\usepackage[colorlinks=true,linkcolor=blue!50!black,citecolor=blue!50!black,urlcolor=blue!50!black]{hyperref}

\graphicspath{{../../}{../../results/}{../../results/topology/}}

% ---- shorthand ---------------------------------------------------------------
\newcommand{\Tr}{\operatorname{Tr}}
\newcommand{\su}{\mathfrak{su}}
\newcommand{\qhat}{\hat{q}}
\newcommand{\Qhat}{\hat{Q}}
\newcommand{\qflow}{q_{t}}
\newcommand{\Qflow}{Q_{t}}
\newcommand{\TA}{\mathrm{TA}}
\newcommand{\Lball}{\mathcal{B}_R}
\newcommand{\dx}{\Delta x}

\title{The Flow-Free Topological Charge Density:\\
A Task Designed to Separate Gauge-Equivariant Attention\\
from Gauge-Equivariant Convolution}
\author{Francesco Passante}
\date{September 14, 2026}

\begin{document}
\maketitle

\begin{abstract}
The matched-parameter shootout on the $0^{++}$ glueball closed on parity: a
gauge-equivariant convolutional network (L-CNN) and the gauge-equivariant
attention network (GELT) produce the same variational operator to within
$1.1\sigma$, at the same mass, with the same signal-to-noise, and the L-CNN does
it for a quarter of the cost. This document argues that the outcome was
structural rather than accidental, extracts from it a selection rule for tasks on
which attention \emph{can} separate, and specifies one such task in full: the
prediction of the \emph{gradient-flowed topological charge density} $\qflow(x)$
from \emph{unflowed} gauge links. The task is per-site rather than
zero-momentum projected, its targets are sparse objects whose position and size
vary from configuration to configuration, and its central difficulty --- telling
a genuine instanton from a lattice dislocation --- is a comparison of field
strength \emph{across scales at one site}, i.e.\ a ratio rather than a
difference. Softmax attention computes ratios across offsets natively; a gated
convolution computes differences and must calibrate them to an amplitude that
varies by orders of magnitude across an ensemble. We give the physics, the
learning problem, the exact architectural statement of what differs between the
two arms, the experimental design with readings fixed in advance, the
implementation plan against the existing codebase, and an honest accounting of
the ways this can fail. The architecture comparison is a coin flip at roughly
$50$--$55\%$; the physics deliverable --- whether a flow-free learned
topological density works at all, and how much of flowed topology is determined
by a bounded neighbourhood --- does not depend on which arm wins.
\end{abstract}

\tableofcontents

% ==============================================================================
\section{Why the glueball shootout tied}
% ==============================================================================

\subsection{What was measured}

The shootout (\texttt{notes/lcnn\_shootout.md} \S9) trained a matched-parameter
L-CNN on the identical problem as GELT: same ensemble, same contiguous
chain-ordered split, same multi-level smeared input channels, same multi-$\Delta$
Rayleigh loss with the same scale pin, same validation-loss checkpoint
selection, same blocked jackknife. The match is on four axes --- parameters
(real degrees of freedom within $10\%$), receptive field (Manhattan $8$ either
way), loop degree ($\le 16$ after four layers either way), and inputs --- and the
one asymmetry left standing is the thing under test: L-Conv's steps are
axis-aligned, so the L-CNN reaches off-axis sites only by stacking, whereas
GELT's shortest-path-averaged transport reaches the whole $L^1$ ball inside one
layer.

The result, combining two independently sampled ensembles of $400$ held-out
configurations each:
\begin{center}
\begin{tabular}{lccc}
\toprule
 & ens0 & ens1 & combined \\
\midrule
$\Delta A_0(\text{GELT} - \text{GEVP})$   & $+0.066 \pm 0.031$ & $+0.089 \pm 0.030$ & $+0.077 \pm 0.022$ \\
$\Delta A_0(\text{L-CNN} - \text{GEVP})$  & $+0.062 \pm 0.031$ & $+0.072 \pm 0.026$ & $+0.068 \pm 0.020$ \\
$\Delta A_0(\text{GELT} - \text{L-CNN})$  & $+0.004 \pm 0.008$ & $+0.016 \pm 0.013$ & $\mathbf{+0.007 \pm 0.007}$ \\
\bottomrule
\end{tabular}
\end{center}
with $\Delta m$ consistent with zero to four decimal places. Both learned
operators beat the classical multi-level GEVP by the same margin; neither beats
the other.

Parity extends further than $A_0$. Measuring the blocked-jackknife
noise-to-signal ratio $\delta C(\Delta)/C(\Delta)$ of the two learned operators
on the \emph{same} $400$ configurations gives
\begin{center}
\begin{tabular}{lcccc}
\toprule
arm & $\Delta = 1$ & $\Delta = 2$ & $\Delta = 4$ & $\Delta = 6$ \\
\midrule
GELT  & $0.0322$ & $0.0396$ & $0.0720$ & $0.1695$ \\
L-CNN & $0.0322$ & $0.0399$ & $0.0728$ & $0.1676$ \\
\bottomrule
\end{tabular}
\end{center}
--- indistinguishable. This is worth recording because it corrects a natural
inference from \S9.2 of the shootout note: the observation there that three of
nine L-CNN operators concentrate $36$--$98\%$ of $C(0)$ on a single
configuration, where zero of fourteen GELT operators exceed $1.0\%$, is a
statement about \emph{failure rate across training runs}, not about typical
variance. It is a reliability argument. It is not an error-bar argument, and no
proposal should be built on the claim that attention gives smaller errors on
this task.

\subsection{The diagnosis}

Three properties of the $0^{++}$ task each independently neutralise attention.

\paragraph{(i) The loss is zero-momentum projected.} The Rayleigh quotient
$-C(1)/C(0)$ is built from the lattice-summed operator
$\bar{O}(t) = \sum_{\vec{x}} O(t,\vec{x})$. Summing over sites is precisely the
operation that discards site-to-site variation, which is the one thing an
input-dependent kernel has and a fixed kernel does not.
\texttt{notes/attention\_as\_operator.md} \S1 made this argument to explain why
the attention-range statistic $\ell_{\mathrm{att}}$ failed twice --- it averaged
$\alpha$ over the lattice and therefore measured the learned \emph{kernel}, a
design quantity of the same kind as a convolution's filter width. The identical
argument applies to the operator itself. Whatever the attention does
differently at site $\vec{x}$ is integrated away before the loss sees it.

\paragraph{(ii) The target is a linear functional of Wilson loops.} The operator
decomposition (\texttt{notes/operator\_decomposition.md};
\texttt{notes/fable5.1\_10-09\_audit.md} \S8.3) puts only $2.5\%$ of
$\|O_{\mathrm{GELT}}\|^2$ outside the span of the $21$-operator classical basis,
and roughly $80\%$ of that residual is rectangular loop shapes the network
rediscovered. The optimum is a fixed linear combination of loop shapes with
configuration-independent coefficients. That is exactly what L-Conv is.

\paragraph{(iii) One coupling, one scale, the most symmetric channel.} The
$0^{++}$ ground state is the most featureless state in the theory, measured at a
single $\beta$ on a single smearing ladder. A single fixed kernel is optimal by
construction; there is nothing to condition on.

Under (i)+(ii)+(iii), softmax attention degenerates into a convolution with
learned per-offset weights --- which is L-Conv. Parity was guaranteed before the
GPU started, and the pre-registration in \texttt{reports/curve} \S7 said so.

\subsection{The selection rule}

The architectural difference that survives is narrow and exact. Write a GELT
block's value path and an L-CNN block's convolution side by side:
\begin{align}
  \text{GEMHSA:}\quad & W'(x) \;=\; \sum_{\dx \in \Lball}
      \alpha_{\dx}(U;x)\; Q_v^{\dagger}(x)\, T_{\dx}(x)\, V(x+\dx)\, T_{\dx}^{\dagger}(x),
      \qquad \alpha = \operatorname{softmax}_{\dx}\!\Big(\tfrac{\operatorname{Re}\Tr[Q^{\dagger}\tilde{K}]}{\sqrt{d_{qkv}\,n_c}}\Big),
      \label{eq:gemhsa} \\
  \text{L-Conv:}\quad & W'(x) \;=\; \sum_{\mu,k}
      w_{\mu k}\; U^{(k)}_{\mu}(x)\, W(x+k\hat{\mu})\, U^{(k)\dagger}_{\mu}(x),
      \qquad w \ \text{constant.}
      \label{eq:lconv}
\end{align}
It is tempting --- and wrong --- to summarise this as ``L-CNN is a polynomial in
the loops, GELT is not''. The L-Act gate
\begin{equation}
  W_{c} \;\longmapsto\; g\!\left(\tfrac{1}{n_c}\operatorname{Re}\Tr W_{c}\right) W_{c},
  \qquad g \in \{\mathrm{ReLU},\ \mathrm{softplus}\},
  \label{eq:lact}
\end{equation}
is itself \emph{input-dependent}, and because L-Conv mixes channels before it,
the gate's argument is a learned linear combination of transported channels. An
L-CNN can therefore realise gates of the form
$\mathrm{ReLU}\big(\operatorname{Re}\Tr[A - B]\big)$, i.e.\ input-dependent
\emph{difference} tests between two scales held in different channels. And GELT
uses the same L-Act gate on its residual branch. So:

\begin{quote}
\textbf{The only architectural difference between the two arms is whether the
weights over \emph{offsets} are input-dependent.} Gating, bilinearity,
receptive field, loop degree, inputs, loss, splits and estimator are matched.
\end{quote}

That is a cleaner experiment than the loose version, and it narrows the claim to
one mechanism. The mechanism is \emph{normalisation}. The softmax in
\eqref{eq:gemhsa} divides by $\sum_{\dx} e^{s_{\dx}}$, so $\alpha$ depends only
on score \emph{differences} and is invariant under a common additive shift of
all scores at a site: it is a relative, scale-free comparison across the
neighbourhood, recomputed per site per configuration. The L-Act gate
\eqref{eq:lact} is not: $\mathrm{ReLU}(\operatorname{Re}\Tr[A-B])$ fires or not
depending on the absolute magnitude of $A - B$, so any decision it implements
must be calibrated to an amplitude that the fixed weights cannot adapt.

Hence the selection rule:
\begin{quote}
\emph{Attention can separate from convolution when the optimal neighbourhood
weighting is \textbf{relative} rather than absolute --- when the right answer
depends on the ratio of local structures rather than on their magnitudes, and
that ratio varies from site to site and from configuration to configuration.}
\end{quote}
Equivalently, a candidate task should have three properties, none of which
$0^{++}$ has: a \textbf{per-site or conditional target} (not lattice-summed);
\textbf{sparse structures at configuration-dependent positions and scales}, so
the useful support moves and resizes; and an \textbf{amplitude-invariant
identity}, so that a fixed-coefficient decision rule must be calibrated and a
normalised one need not.

There is already one measurement in this repository consistent with that rule.
At $\beta = 0.760$ in $3$D $\mathbb{Z}_2$, the coupling closest to $\beta_c$ in
the scan, the classical fixed-radius smeared operator has essentially no signal
($C(1)/C(0) \approx 0.007$) and an untrained network's attention field has none
either, while the \emph{trained} attention field gives a clean exponential
($0.496,\,0.394,\,0.329,\,0.276$) --- see
\texttt{notes/attention\_as\_operator.md} \S6.1. That is the regime where a
fixed kernel stops working and a learned relative one does not. It is
suggestive, not decisive: it concerns the attention \emph{field} rather than an
output operator, and no L-CNN arm was ever run there.

% ==============================================================================
\section{Topology on the lattice}
% ==============================================================================

\subsection{The continuum object}

For a smooth $SU(N)$ gauge field on a compact four-manifold the topological
charge
\begin{equation}
  Q \;=\; \frac{1}{32\pi^2}\int d^4x\; \epsilon_{\mu\nu\rho\sigma}
          \operatorname{Tr}\!\left[F_{\mu\nu}(x) F_{\rho\sigma}(x)\right]
  \;\in\; \mathbb{Z}
  \label{eq:Qcont}
\end{equation}
is an integer, and the integrand $q(x)$ --- the topological charge density --- is
a local, gauge-invariant, parity-odd scalar field. Its physics content is
substantial: the topological susceptibility $\chi_{\mathrm{top}} = \langle Q^2
\rangle / V$ enters the Witten--Veneziano relation for the $\eta'$ mass, the
$\theta$-dependence of the vacuum energy is controlled by the same distribution,
and the two-point function $\langle q(x) q(0) \rangle$ has the celebrated
structure of a positive contact term plus a negative-definite tail, which is
what makes $\chi_{\mathrm{top}} > 0$ compatible with reflection positivity.

\subsection{Lattice discretisations, and one that is in this repository}

The Wilson gauge action has no exact topology: the configuration space is
connected, and any lattice definition of $Q$ is a real number that approaches an
integer only on sufficiently smooth fields. The simplest definition replaces
$F_{\mu\nu}$ by the anti-Hermitian part of the plaquette,
\begin{equation}
  F^{\mathrm{plaq}}_{\mu\nu}(x) \;=\; \frac{P_{\mu\nu}(x) - P^{\dagger}_{\mu\nu}(x)}{2i},
  \qquad
  q^{\mathrm{plaq}}_x \;=\; \frac{\epsilon_{\mu\nu\rho\sigma}}{32\pi^2}
    \operatorname{Tr}\!\left[F^{\mathrm{plaq}}_{\mu\nu}(x) F^{\mathrm{plaq}}_{\rho\sigma}(x)\right],
  \label{eq:qplaq}
\end{equation}
which is exactly \texttt{gelt.lattice.topological\_charge\_density} and exactly
Eq.~(13) of the L-CNN paper. It carries $O(a^2)$ discretisation errors and a
multiplicative renormalisation, $Q^{\mathrm{latt}} \simeq Z(\beta)\,Q + \text{noise}$
with $Z < 1$ on rough fields.

\paragraph{A defect worth recording.} Definition \eqref{eq:qplaq} bases all six
plaquettes at the \emph{corner} $x$, so it is not symmetric under reflections
about $x$ and is therefore not exactly parity-odd --- not even after summing
over the lattice. We checked this directly. Build the reflection
$x_1 \mapsto (-x_1) \bmod L$ with the correct link relabelling
$U'_{\nu}(y) = U_{\nu}(Py)$ for $\nu \neq 1$ and
$U'_{1}(y) = U_{1}\!\big(P(y+\hat{1})\big)^{\dagger}$. This is an exact symmetry
of the Wilson action: on a $6^4$ $SU(2)$ configuration the action is preserved to
all printed digits, $15478.81704264 \to 15478.81704264$. Under the same map the
charge \eqref{eq:qplaq} goes $-1.3161 \to +2.9888$: the sign flips, the magnitude
does not. The remedy is the standard \textbf{clover} discretisation, which
averages the four plaquettes in each plane around $x$ so that $F_{\mu\nu}$ is
centred on the site,
\begin{equation}
  F^{\mathrm{clov}}_{\mu\nu}(x) \;=\; \frac{1}{8i}
    \sum_{\text{4 leaves}} \left(P^{(i)}_{\mu\nu}(x) - P^{(i)\dagger}_{\mu\nu}(x)\right),
  \label{eq:qclov}
\end{equation}
restoring both the hypercubic symmetry and the exact parity-oddness of $q_x$ and
of $Q$. It does not exist in the repository and is a prerequisite for this
study --- both for the target and for the sanity checks. It is also the natural
home of a new test in \texttt{tests/test\_lattice.py}: \emph{the clover charge is
parity-odd to machine precision and the plaquette charge is not}.

\subsection{Dislocations, and why the naive density is not the object of interest}

The real problem is not discretisation error but \emph{dislocations}: localised
lattice artifacts of size of order $a$ on which $|q^{\mathrm{plaq}}_x|$ is large
while nothing topological is present. For the Wilson action the lattice action of
an instanton of radius $\rho$ falls below the continuum bound $8\pi^2/g^2$ once
$\rho \lesssim a$, so such objects are not Boltzmann-suppressed and contribute a
spurious ultraviolet piece to $\chi_{\mathrm{top}}$ that does not scale. On a
thermalised configuration, $q^{\mathrm{plaq}}$ is dominated by them.

This has two consequences that matter here. First, the physically meaningful
object is not $q^{\mathrm{plaq}}$ but the density of a \emph{smoothed}
configuration. Second --- and this is the whole point of the present proposal ---
the operation that separates a genuine instanton from a dislocation is
intrinsically a comparison \emph{across scales at one site}: is the lump extended
over several lattice spacings, or is it a single-plaquette spike? A large
$|q_x|$ alone says nothing. The ratio of the field strength averaged at radius
$a$, $2a$, $3a$ says everything. That is a scale-free test, and its threshold
cannot be an absolute magnitude, because $|q^{\mathrm{plaq}}_x|$ spans orders of
magnitude across an ensemble and across $\beta$.

\subsection{The gradient flow}

The standard smoother is the Wilson (gradient) flow,
\begin{equation}
  \dot{V}_{t}(x,\mu) \;=\; -g_0^2\,\partial_{x,\mu} S_W\!\left[V_t\right]\; V_t(x,\mu),
  \qquad V_0 = U,
  \label{eq:flow}
\end{equation}
a deterministic diffusion whose effect at flow time $t$ is to average the field
over a physical radius
\begin{equation}
  r_{\mathrm{sm}} \;=\; \sqrt{8t}.
  \label{eq:rsm}
\end{equation}
On the lattice the drift is read straight off the staple. With
$A_{\mu}(x)$ the staple sum in the repository's convention --- so that the local
action is $S_{\mathrm{local}} = -(\beta/n_c)\operatorname{Re}\Tr[U_{\mu}(x) A_{\mu}(x)]$ ---
one integrates
\begin{equation}
  U_{\mu}(x) \;\longleftarrow\; \exp\!\big(\varepsilon\, Z_{\mu}(x)\big)\, U_{\mu}(x),
  \qquad
  Z_{\mu}(x) \;=\; -\Big[\,U_{\mu}(x)\,A_{\mu}(x)\,\Big]_{\TA},
  \label{eq:flowstep}
\end{equation}
where $[M]_{\TA} = \tfrac{1}{2}(M - M^{\dagger}) - \tfrac{1}{2N}\Tr(M - M^{\dagger})\mathbb{1}$
is the traceless anti-Hermitian projection. For $n_c = 2$ both $[\cdot]_{\TA}$
and the exponential are closed-form --- $\exp(i\theta\, \hat{n}\cdot\vec{\sigma})
= \cos\theta + i \sin\theta\, \hat{n}\cdot\vec{\sigma}$ --- so no \texttt{linalg}
kernel is needed, in keeping with the repository's existing
$SU(2)$ fast paths.

The flow is \emph{not} a linear smoothing. In the small-field regime
\eqref{eq:flow} linearises to a heat kernel of radius $\sqrt{8t}$, but on
thermalised fields at $\beta \sim 2.4$, where the mean plaquette is far from
unity, it is strongly nonlinear, and the nonlinearity is exactly what removes
dislocations: a sub-lattice-size instanton shrinks under the flow and falls
through the lattice, while a genuine one of radius $\rho \gg a$ is preserved.
That asymmetry --- destroy the spike, keep the lump --- is the part of the map
that a linear filter cannot represent, and it is the part on which this
experiment rests.

\subsection{Why the \emph{density} and not just \texorpdfstring{$Q$}{Q}}

Predicting the single integer $\Qflow$ is a substantially easier and largely
solved problem (\S\ref{sec:prior}). The field $\qflow(x)$ is a strictly richer
object and it is what the physics wants: the instanton size distribution, which
at one loop behaves as $n(\rho)\,d\rho \propto \rho^{\,b-5} d\rho$ with
$b = 11N/3$ (so $\rho^{7/3}$ for $SU(2)$) until infrared effects cut it off; the
local structure of the vacuum (instanton-liquid versus sign-coherent sheets); and
the two-point function $\langle q(x) q(0)\rangle$, whose contact term and
negative tail are the subject of a long literature and which is badly distorted
by the very smoothing used to define it. A learned map from $\tau = 0$ gives, in
principle, a density at the original lattice spacing --- something flow cannot
provide by construction.

% ==============================================================================
\section{The learning problem}
% ==============================================================================

\subsection{Definition}

Fix a gauge group $SU(2)$, a coupling $\beta$, a lattice $\Lambda = L^4$, and a
flow time $t$. Let $\mathcal{F}_t$ denote the Wilson flow \eqref{eq:flow}
integrated to $t$. The task is to learn
\begin{equation}
  \qhat \;:\; U \;\longmapsto\; \qhat(x)[U] \;\approx\; \qflow(x)
  \;\equiv\; q^{\mathrm{clov}}_x\!\left[\mathcal{F}_t(U)\right],
  \qquad U \ \text{a thermalised configuration at } \tau = 0,
  \label{eq:task}
\end{equation}
by minimising a per-site mean-squared error over the lattice and the ensemble.
The prediction of the total charge follows as $\Qhat = \sum_{x} \qhat(x)$.

Three things about \eqref{eq:task} deserve emphasis. The input is
\textbf{unflowed} --- the raw Monte Carlo configuration. The target is the
\textbf{density field}, not a scalar. And the map is a genuine inference: the
network must perform, implicitly and in one shot, the work that the flow does
over hundreds of integrator steps.

\subsection{Is it well posed?}

Topological charge is a global invariant, so it is worth being explicit about why
a \emph{local} network can learn $\qflow(x)$ at all. The flow is a parabolic
equation; its Green function is exponentially localised on the scale
$r_{\mathrm{sm}} = \sqrt{8t}$. Consequently $\qflow(x)$ depends on the unflowed
links in a neighbourhood of $x$ of radius $O(\sqrt{8t})$ up to exponentially
small tails. The task is therefore well posed as a local regression \emph{provided}
\begin{equation}
  \sqrt{8t}\,/\,a \;\lesssim\; R_{\mathrm{net}},
  \label{eq:wellposed}
\end{equation}
where $R_{\mathrm{net}}$ is the network's receptive field --- Manhattan $8$ for
both arms at $R = 2$, $K = 2$ and four layers. Condition \eqref{eq:wellposed} is
not a nuisance; it is a \emph{measurement}. Sweeping $t$ and watching where each
architecture's accuracy collapses answers the physics question ``how much of
flowed topology is determined by a bounded neighbourhood of the unflowed
field?'', and doubles as a capacity control (\S\ref{sec:readings}).

Note also that no network can be expected to reproduce the integer quantisation
of $\Qflow$ exactly, since that is a global property emerging from a sum of local
predictions. How close $\Qhat$ comes to an integer is a diagnostic of quality,
not a constraint that is imposed.

\subsection{What this is not}
\label{sec:notthis}

It is important to distinguish \eqref{eq:task} from the topology task in the
L-CNN literature, because the names collide. Favoni et al.\ regress
$q^{\mathrm{plaq}}_x$ \emph{of the configuration the network is shown}. That is
a local algebraic identity: \eqref{eq:qplaq} is a degree-two polynomial in
plaquettes at a single site, and a single L-Bilin layer spans it exactly. Their
demonstration applies the trained model \emph{along} a Wilson flow trajectory ---
at each flow time it is fed the flowed links and asked for the flowed naive
density --- and the integer values in their figure come from the flow, not from
the network. The map is input-to-same-configuration. Ours is
input-at-$\tau=0$-to-target-at-$\tau=t$: nonlocal, nonlinear, and
scale-selective. The L-CNN's proven competence at the former says very little
about the latter, which is precisely why the comparison is worth running.

% ==============================================================================
\section{Why this task can discriminate}
% ==============================================================================

Against the selection rule of \S1.3, the task scores on all three properties.

\paragraph{Per-site target.} The loss is $\sum_x (\qhat(x) - \qflow(x))^2$.
Nothing is summed over the lattice before the gradient is taken, so the
site-to-site variation that the $0^{++}$ Rayleigh quotient integrated away is
exactly what is being supervised. Tie-cause (i) is removed by construction.

\paragraph{Sparse structures at configuration-dependent positions and scales.}
Instantons are dilute, localised, and distributed in size. Within a single
ensemble $\rho/a$ already spans a range, and across the $\beta$ ladder the mean
$\rho/a$ moves by roughly a factor of two. The neighbourhood that the network
must aggregate over is therefore not fixed: it is $\rho(x)$, a quantity that
varies within a configuration and across an ensemble. Tie-cause (iii) is removed.

\paragraph{Amplitude-invariant identity.} This is the load-bearing one. The
decision ``is this lump topological or is it a dislocation?'' is invariant under
rescaling the local field strength --- it depends on the \emph{shape} of the
profile, i.e.\ on ratios of field strength at radius $a$, $2a$, $3a$, not on
their absolute size. Now compare the two architectures' available decision rules
at one site:
\begin{align}
  \text{GELT:}\quad & \alpha_{\dx} = \frac{e^{s_{\dx}}}{\sum_{\dx'} e^{s_{\dx'}}}
    \qquad\text{---\ invariant under } s \to s + c, \ \text{a pure ratio across offsets;} \\
  \text{L-CNN:}\quad & g\!\left(\operatorname{Re}\Tr\!\left[\textstyle\sum_{c} w_c\, W_c\right]\right)
    \qquad\text{---\ a threshold on an absolute magnitude.}
\end{align}
An L-CNN can hold a radius-$a$ average in one channel and a radius-$3a$ average
in another, and it can gate on their difference. What it cannot do with fixed
weights is gate on their \emph{ratio}: the threshold in
$\mathrm{ReLU}(\operatorname{Re}\Tr[A - B])$ is an absolute number, so the
decision is correct only for configurations whose local field strength sits near
the scale the weights were calibrated to. Since $|q^{\mathrm{plaq}}_x|$ is
heavy-tailed within an ensemble and shifts systematically with $\beta$, that
calibration cannot be right everywhere.

This is a statement about efficiency, not about representability. With enough
channels and depth an L-CNN can approximate a ratio on any bounded domain; the
claim is that at a matched budget of $\sim 15$--$26$k real degrees of freedom it
will spend capacity on a filter bank plus a calibrated selector where GELT gets
the selector for free from the softmax. That is a bet, and \S\ref{sec:risks}
prices it honestly.

% ==============================================================================
\section{Prior art, and what is open}
\label{sec:prior}
% ==============================================================================

\begin{center}
\begin{tabular}{p{0.30\textwidth} p{0.28\textwidth} p{0.34\textwidth}}
\toprule
work & input $\to$ target & status \\
\midrule
Favoni, Ipp, M\"uller, Schuh, \emph{PRL} \textbf{128} (2022) 032003
 & flowed links $\to$ $q^{\mathrm{plaq}}$ of \emph{the same} configuration
 & solved; local, degree-two, algebraic. One L-Bilin layer spans it. \\[2pt]
Favoni, Ipp, M\"uller (2022), \texttt{arXiv:2212.00832}
 & links $\to$ links, neural-ODE Wilson flow
 & toy models only; a different problem (learning the flow map itself). \\[2pt]
Matsumoto et al., \emph{PTEP} \textbf{2021} 023D01 ($SU(3)$)
 & $q$-density at \emph{small} flow time $t/a^2 \le 0.3$ $\to$ integer $Q$ at large $t$
 & $>99\%$; scalar target, partial flow, non-equivariant networks. Best results integrate out space. \\[2pt]
\midrule
\textbf{this proposal}
 & \textbf{unflowed links ($\tau = 0$) $\to$ flowed density field $\qflow(x)$}
 & \textbf{open.} \\
\bottomrule
\end{tabular}
\end{center}

Two honest consequences follow. First, the gap is real: nobody has learned the
flowed \emph{density field} from a genuinely unflowed configuration, and the
L-CNN's published topology success is on the algebraic task, not this one.
Second, the \emph{compute-saving} motivation is weaker than it first appears.
Matsumoto et al.\ show that a little flow plus a simple classifier already
recovers the integer $Q$ with high accuracy, so ``skip the expensive flow'' is
partly taken. The motivation that survives is the density field itself
(\S2.5) and the locality question \eqref{eq:wellposed}. The write-up must lead
with those, not with a cost argument it cannot fully support.

% ==============================================================================
\section{Experimental design}
% ==============================================================================

\subsection{Ensembles and scale}

Three couplings, with the lattice size chosen so that the physical volume is
approximately constant, exploiting the fact that both architectures are
translation-equivariant and can be trained and evaluated at any size:

\begin{center}
\begin{tabular}{cccccc}
\toprule
$\beta$ & $a\sqrt{\sigma}$ (lit.) & $a$ [fm] & $L$ & $La$ [fm] & configurations \\
\midrule
$2.3$ & $\approx 0.369$ & $\approx 0.165$ & $8$  & $1.32$ & $1200$ \\
$2.4$ & $\approx 0.266$ & $\approx 0.119$ & $12$ & $1.43$ & $1200$ \\
$2.5$ & $\approx 0.188$ & $\approx 0.084$ & $16$ & $1.35$ & $1200$ \\
\bottomrule
\end{tabular}
\end{center}

The $a\sqrt{\sigma}$ values are the standard Wilson-action $SU(2)$ numbers and
are quoted here only to fix the design; \textbf{they must be re-derived in-repo}
from \texttt{rectangular\_wilson\_loop} before any physical statement is made, in
keeping with the repository's practice of never importing a scale it has not
measured. The lattice spacing changes by a factor of two across the ladder, so
$\rho/a$ --- the object size in the units the network sees --- changes by the
same factor. That is the scale-adaptivity axis.

Sampling uses the existing exact $SU(2)$ heat-bath plus overrelaxation
(\texttt{heatbath\_overrelaxation\_sweep}), which is the prerequisite for
decorrelated topology: topological modes have the longest autocorrelation time in
the theory, and $n_{\mathrm{skip}}$ must be set from
\texttt{integrated\_autocorrelation\_time} applied to $Q$ itself, not to the
plaquette. \textbf{Topological freezing is a live risk at $\beta = 2.5$} and must
be measured, not assumed: if $\tau_{\mathrm{int}}(Q)$ is large the effective
statistics collapse and the $\beta = 2.5$ row becomes a single-sector dataset,
which would invalidate the aggregate readings while leaving the per-site ones
intact.

\subsection{Flow, and the target ladder}

The flow \eqref{eq:flowstep} is integrated at fixed \textbf{lattice-unit} flow
time $t/a^2$, identical across $\beta$. This is deliberate and it is the design
decision that makes the experiment interpretable. Fixing $t/a^2$ holds
$r_{\mathrm{sm}}/a$ --- and therefore the receptive-field requirement
\eqref{eq:wellposed} --- constant across the ladder, so that any degradation of a
mixed-$\beta$ model is attributable to \emph{scale adaptivity} and not to a
capacity wall that bites at one $\beta$ and not another. Fixing $t$ physically
would confound the two.

The primary target is $t/a^2 = 2$, i.e.\ $r_{\mathrm{sm}}/a = 4$, comfortably
inside Manhattan $8$. The full ladder
\begin{equation}
  t/a^2 \;\in\; \{0.5,\ 1,\ 2,\ 4,\ 8\},
  \qquad r_{\mathrm{sm}}/a \;=\; \{2.0,\ 2.8,\ 4.0,\ 5.7,\ 8.0\}
\end{equation}
is generated once per configuration and stored, because it costs one flow
trajectory and yields five targets. The largest rung deliberately sits at the
receptive-field boundary.

\paragraph{Gate 0.} Before any training, verify the $Q(t)$ plateau at every
$\beta$: flow to $t/a^2 = 16$, plot $Q(t)$, and confirm that $t/a^2 = 2$ lies on
a plateau and that $Q$ is within a stated tolerance of an integer. If it does
not at $\beta = 2.3$ --- which is plausible, the coarsest lattice being the one
where dislocations survive longest --- the primary rung moves, and the fact is
recorded rather than hidden.

\subsection{Splits}

Contiguous and chain-ordered, matching the glueball protocol: $800$ train,
$200$ validation, $200$ test, with the test block at the far end of the Markov
chain. This is deliberately the harder split; \texttt{notes/lcnn\_shootout.md}
\S9.1 shows why it matters (a network can fit chain-local structure that survives
into an adjacent validation block and not into a distant test block). Per-site
supervision gives $800 \times L^4$ training targets, so the statistics per
configuration are enormous and small batches suffice.

\subsection{Matched budgets}

The four-axis match of \texttt{notes/lcnn\_shootout.md} \S2 carries over, with
$D = 4$ substituted:

\begin{center}
\begin{tabular}{lll}
\toprule
axis & GELT & L-CNN \\
\midrule
receptive field & $R = 2$ $L^1$ ball $\times$ 4 layers $\to$ Manhattan 8 & $K = 2$ symmetric $\times$ 4 layers $\to$ Manhattan 8 \\
loop degree & bilinear value path, $\le 16$ & L-CB doubling, $\le 16$ \\
inputs & $6$ plaquette channels ($n_{\mathrm{pairs}} = D(D{-}1)/2$) & identical \\
parameters & tune $d_{\mathrm{model}}$ & tune $c_{\mathrm{hidden}}$ to within $15\%$ \\
output & \texttt{reduction="none"} $\to (B, \Lambda)$ & \texttt{reduction="none"} $\to (B,\Lambda)$ \\
\bottomrule
\end{tabular}
\end{center}

Two $D = 4$ specifics. The $L^1$ ball at $R = 2$ in four dimensions has
$40$ offsets against the glueball task's $24$ in three, so the transport stage ---
already $62.8\%$ of a GELT step
(\texttt{notes/performance\_audit.md} \S5.0) --- grows by $1.67\times$ per site.
And \textbf{RoPE requires $d_{qkv} \ge 2D = 8$} (known caveat 2): below that,
whole lattice axes sit on the identity rotation and the score stops being
offset-selective along them. This is a hard constraint on the GELT arm's
geometry and must be asserted in the constructor, not remembered.

\subsection{Arms}

Six arms, of which two are classical and cost nothing:

\begin{enumerate}
\item \textbf{Identity}: $\qhat = q^{\mathrm{clov}}_x[U]$, the unflowed density.
      The zero-work baseline; everything must beat it.
\item \textbf{Best linear filter}: the least-squares-optimal isotropic linear
      filter of support Manhattan $8$ applied to $q^{\mathrm{clov}}[U]$, fitted on
      the training split. This is the ``a convolution is the right answer''
      hypothesis made concrete and is the single most important control in the
      design: if the networks do not clearly beat it, nothing nonlinear has been
      learned and the architecture comparison is moot.
\item \textbf{GELT, per-$\beta$} (three models).
\item \textbf{L-CNN, per-$\beta$} (three models).
\item \textbf{GELT, mixed-$\beta$} (one model trained on all three ensembles).
\item \textbf{L-CNN, mixed-$\beta$} (one model).
\end{enumerate}
Plus untrained controls of both architectures at three seeds each, following the
precedent of \texttt{notes/fable5.1\_10-09\_audit.md} \S8.3, with the
seed-aggregation rule (median, not mean) fixed \emph{before} the numbers are
seen --- \S9.2 of the shootout note records why.

\subsection{Metrics}

\begin{description}
\item[M1 --- per-site accuracy.] $1 - R^2$ on the test split, i.e.\
  $\sum_x (\qhat - \qflow)^2 / \sum_x (\qflow - \langle \qflow\rangle)^2$.
  Blocked jackknife over configurations, block size from $\tau_{\mathrm{int}}(Q)$.
\item[M2 --- charge accuracy.] $\mathrm{RMS}(\Qhat - \Qflow)$ and the RMS
  distance of $\Qhat$ to the nearest integer.
\item[M3 --- physics.] $\chi_{\mathrm{top}}$ and the histogram of $Q$
  reconstructed from $\qhat$ versus from $\qflow$, on the same configurations.
\item[M4 --- stratified error.] M1 restricted to the dislocation-rich tail:
  sites in the top percentile of $|q^{\mathrm{clov}}[U]|$, and configurations in
  the top decile of roughness. \emph{This is where the mechanism predicts the
  separation, and it is the reading to look at first.}
\item[M5 --- the two-point function.] $\langle \qhat(x)\qhat(0)\rangle$ against
  $\langle \qflow(x)\qflow(0)\rangle$, including the sign of the tail.
\item[M6 --- robustness.] Top-1 configuration share of the loss, under the
  uniform gate installed in \texttt{fit\_glueball\_overlap.py} by
  \texttt{notes/lcnn\_shootout.md} \S9.2.
\end{description}

All GELT-versus-L-CNN differences are taken \emph{inside} each jackknife sample
on shared configurations, as \texttt{fit\_glueball\_overlap.py --vs} does: the
two arms are correlated through the ensemble, and the correlated error measured
roughly $4\times$ smaller than the naive one on the glueball task.

\subsection{Readings, fixed in advance}
\label{sec:readings}

\begin{description}
\item[R1 (primary, architecture).] $\Delta(1-R^2)$ between mixed-$\beta$ GELT and
  mixed-$\beta$ L-CNN, correlated jackknife. \textbf{GELT wins} if the relative
  error is lower by more than $2\sigma$. \textbf{Parity} if within $1\sigma$.
\item[R2 (control, adaptivity).] The same difference at fixed $\beta$.
  \emph{Prediction: parity, or a margin clearly smaller than R1's.} If single-$\beta$
  already separates, the effect is not scale adaptivity and the claim must be
  re-stated --- possibly in GELT's favour, but for a different reason.
\item[R3 (control, capacity).] Accuracy across the flow-time ladder. Both arms
  must collapse as $r_{\mathrm{sm}}/a \to 8$. If GELT's collapse sets in later,
  that is receptive-field efficiency, not adaptivity, and must be reported as
  such.
\item[R4 (physics, architecture-independent).] Does \emph{either} network
  reproduce $\chi_{\mathrm{top}}$ and the $Q$-histogram from $\tau = 0$ within the
  statistical error of the flowed reference? This is the deliverable that does
  not depend on who wins.
\item[R5 (nulls).] The untrained arms must fail. The identity arm and the
  best-linear-filter arm must be clearly beaten. If arm 2 is not beaten, stop:
  the task is a convolution and the paper says so.
\end{description}

None of the five outcomes is a wasted run, which is the argument for spending the
time --- the same argument \S7 of the shootout note made, and it held.

% ==============================================================================
\section{Implementation}
% ==============================================================================

\subsection{What already exists}

More than one would expect. Verified by direct execution at $D = 4$, $SU(2)$,
$L = 6$, \texttt{complex128}:

\begin{itemize}
\item \texttt{plaquette\_tensor} returns $(B, 6, \Lambda, 2, 2)$ in $D=4$;
  \texttt{topological\_charge\_density} returns $(B, \Lambda)$;
  \texttt{build\_transport\_average} returns $(B, 40, \Lambda, 2, 2)$ at $R=2$;
  \texttt{build\_axis\_transports} returns $(B, 4, 2, \Lambda, 2, 2)$.
\item Both models run and are gauge invariant to machine precision: drift
  $0.0$ (GELT), $1.4\times 10^{-14}$ (L-CNN), against $1.5\times10^{-16}$ for the
  target itself.
\item \textbf{\texttt{LCNN} already accepts \texttt{reduction="none"}} and returns
  $(B, \Lambda)$, so no per-site head has to be written for the baseline.
\item \texttt{staple\_sum(..., batched=True)} supplies the flow drift
  \eqref{eq:flowstep} directly, already vectorised over the configuration axis.
\item \texttt{data.build\_plaquette\_datasets} takes \texttt{target} as a
  callable \texttt{target(configs, group) -> Tensor}, so the flowed density binds
  in with \texttt{functools.partial} with no change to the builder, and \texttt{R}
  makes it emit $(X, T, y)$ triples with the transport precomputed.
\item \texttt{integrated\_autocorrelation\_time} sets $n_{\mathrm{skip}}$.
\end{itemize}

\subsection{What must be written}

\begin{enumerate}
\item \texttt{gelt/flow.py} --- \texttt{wilson\_flow(U, group, t, eps,
  integrator)} implementing \eqref{eq:flowstep}, with the closed-form $SU(2)$
  exponential and the traceless anti-Hermitian projection, batched over
  configurations, plus \texttt{flow\_trajectory} returning the target ladder in
  one pass. L\"uscher's third-order Runge--Kutta is preferred over Euler; Euler
  at small $\varepsilon$ is the fallback and the cross-check.
\item \texttt{topological\_charge\_density(..., definition="clover")} in
  \texttt{gelt/lattice.py}, per \eqref{eq:qclov}. \textbf{Required}, not
  optional: \S2.2 shows the existing plaquette definition is not parity-odd.
\item \texttt{scripts/train\_topology.py} --- the trainer, modelled on
  \texttt{train\_glueball.py}'s single-\texttt{\_build\_model()} discipline so
  that \texttt{TOPO\_ARCH=lcnn} switches the architecture and nothing else, and
  so that a profiler cannot drift from the training path. Per-site MSE,
  validation-loss checkpoint selection, artifact names carrying $\beta$, $L$,
  flow time, architecture, width and seed.
\item \texttt{scripts/measure\_topology.py} --- the classical arms (identity,
  best linear filter), gate 0's $Q(t)$ plateau, $\tau_{\mathrm{int}}(Q)$, and the
  ensemble cache.
\item \texttt{scripts/fit\_topology.py} --- offline: M1--M6 from the dumps, with
  the correlated \texttt{--vs} difference.
\item Tests: clover parity-oddness to machine precision (and the plaquette
  definition failing it); flow gauge covariance; flow preserving the group;
  Euler-versus-RK3 agreement; the $D=4$ parameter match between the two arms;
  the $d_{qkv} \ge 2D$ assertion.
\end{enumerate}

\subsection{Cost}

Estimated from the measured V100 numbers in \texttt{notes/performance\_audit.md}
(GELT $5.04$ s/step at $248\,832$ sites $\times\ 24$ offsets; the matched L-CNN
$1.29$ s/step):

\begin{center}
\begin{tabular}{lll}
\toprule
phase & scale & estimate \\
\midrule
ensembles ($3\beta$, heat-bath + OR) & $8^4$, $12^4$, $16^4$, $1200$ each & $\sim 1.5$ nights \\
flow targets (5-rung ladder) & $3600$ trajectories & $1$--$2$ hours \\
GELT trainings & $3$ per-$\beta$ + $1$ mixed & $\sim 2$ nights \\
L-CNN trainings & $3$ per-$\beta$ + $1$ mixed & $\sim 0.5$ night \\
untrained + classical arms & offline & hours \\
\midrule
\multicolumn{2}{l}{\textbf{total}} & $\sim 4$--$5$ V100 nights, $\sim 1$ build week \\
\bottomrule
\end{tabular}
\end{center}

% ==============================================================================
\section{Risks, and what each outcome means}
\label{sec:risks}
% ==============================================================================

\paragraph{The map is mostly a smoothing.} The strongest objection. In the
small-field limit the flow \emph{is} a linear filter, so a convolution is the
right architecture for the bulk of the map and the adaptive part is a correction.
The size of that correction is an empirical question, and arm 2 (best linear
filter) measures it directly: the gap between arm 2 and the best network
\emph{is} the nonlinear content of the flow. If that gap is small, the whole
study reduces to ``the flow is nearly linear at these couplings'', which is
itself a clean, quotable result --- and it terminates the architecture question
honestly rather than burying it.

\paragraph{The L-CNN approximates the ratio well enough.} At four layers and
five or six channels with L-Act gates it may build a serviceable calibrated
selector. \emph{Probability I assign to a clear GELT win on R1: $50$--$55\%$.}
This is down from the $65\%$ quoted before this audit, for three reasons: the
surviving mechanism is one axis rather than four (\S1.3); the map is mostly a
smoothing; and the base rate in this repository for ``attention will help here''
predictions is currently zero for one.

\paragraph{Topological freezing.} At $\beta = 2.5$ the charge may decorrelate so
slowly that the $1200$ configurations contain a handful of independent
topological samples. This does not affect the per-site readings M1/M4, which have
$L^4$ targets per configuration, but it would invalidate M2/M3 at that coupling.
Mitigation: measure $\tau_{\mathrm{int}}(Q)$ first and report it; if freezing
bites, quote the aggregate physics at $\beta = 2.3$ and $2.4$ only, and say why.

\paragraph{Single-$\beta$ already separates.} The instanton size distribution is
broad within one ensemble, so the adaptivity requirement may not need the
$\beta$ ladder at all. This breaks R2 as a control while strengthening the
overall case. The response is to re-stratify by measured $\rho$ rather than by
$\beta$ and to say plainly that the pre-registered control did not behave as
predicted.

\paragraph{The target is ill-defined at the coarsest coupling.} If $Q(t)$ has no
plateau at $\beta = 2.3$ and $t/a^2 = 2$, the regression target is partly
arbitrary. Gate 0 catches this before any GPU time is spent.

\paragraph{GELT's cost.} The transport is already $62.8\%$ of a GELT step in
three dimensions and grows by $1.67\times$ in four. If the step is intolerable,
\texttt{notes/performance\_audit.md} \S5.1 --- the adjoint $SO(3)$ representation
for $SU(2)$, worth a predicted $\sim 3.5\times$ on the transport and $1.8\times$
end to end --- is the lever, and this task is a better reason to pull it than the
glueball one was.

\paragraph{What each outcome buys.} A GELT win on R1 with parity on R2 is the
designed result: attention buys scale adaptivity, and the thesis has the
architectural claim the $0^{++}$ chapter could not supply. Parity on both is a
second null, and after the shootout a second null is genuinely informative --- it
would say that input-dependent offset weighting does not pay on \emph{local
regression} either, which sharply narrows where the mechanism could ever pay and
is worth writing down. An L-CNN win points at the transport average or the
receptive field rather than at the loss, and should be reported before an
examiner finds it. In all three cases R4 stands on its own.

% ==============================================================================
\section{Summary}
% ==============================================================================

The $0^{++}$ shootout tied because the task had none of the three properties that
let an input-dependent kernel pay: its loss was zero-momentum projected, its
optimum was a linear functional of loop shapes, and it lived at a single scale.
Prediction of the gradient-flowed topological charge density from unflowed links
has all three. It is per-site; its targets are sparse objects whose position and
size vary within and across ensembles; and its central operation --- separating a
genuine instanton from a lattice dislocation --- is a scale-free ratio test,
which is what a softmax over offsets computes natively and what a gated
convolution must calibrate. The architectural difference between the two arms
reduces, exactly, to whether the offset weights depend on the input, which makes
the experiment a controlled A/B on one mechanism.

The honest odds on the architecture question are close to even. The reason to run
it anyway is that the physics deliverable is separable and largely independent of
the outcome: no one has learned the flowed topological charge \emph{density}
from a genuinely unflowed configuration, and the flow-time ladder turns the
receptive-field constraint into a measurement of how much of flowed topology is
determined by a bounded neighbourhood of the raw field. Both are worth having
whichever architecture wins.

\end{document}
